% =============================================================================
% Combined analysis of linezolid potency (EC50_L).
%
% Supersedes and merges two earlier documents:
%   - EC50_lzd_threshold_methods.tex  (deterministic threshold determination,
%     EC50_lzd_threshold_sweep.py)
%   - MC_ec50_lzd_description.tex     (Monte Carlo population, MC_ec50_lzd.py)
%
% All values verified against the code at the current baseline
% (rho_S = 0.72, EC50_L = 0.6, lzd_res_fraction = 0.30).
% =============================================================================
\documentclass[11pt]{article}
\usepackage[margin=1in]{geometry}
\usepackage{amsmath,amssymb}
\usepackage{booktabs}
\usepackage{hyperref}

\title{Analysis of Linezolid Potency ($EC_{50,L}$)\\
\large Deterministic escape threshold and Monte Carlo patient population}
\author{}
\date{}

\begin{document}
\maketitle

\section{Overview}

Linezolid's  $EC_{50,L}$, is analysed
in two complementary ways, implemented in two scripts.

\begin{enumerate}
\item A \textbf{deterministic threshold analysis}
(\texttt{EC50\_lzd\_threshold\_sweep.py}) treats $EC_{50,L}$ as a
bifurcation parameter and locates, by root-finding, the exact value
separating suppression from resistant escape. It answers: \emph{where is
the switch?}
\item A \textbf{Monte Carlo population analysis}
(\texttt{MC\_ec50\_lzd.py}) draws $1000$ virtual patients from a log-normal
$EC_{50,L}$ distribution and simulates each one. It answers: \emph{given
realistic inter-patient variability, what fraction of patients fall on
each side of that switch, and what do their bacterial kinetics look
like?}
\end{enumerate}

The two are independent: the Monte Carlo script derives the threshold
internally, over a different bracketing interval from the threshold analysis. The deterministic threshold analysis calculates the bifurcation point using the bisection method.  Both methods agree, providing a consistency check as discussed in \S\ref{sec:agreement}.

Both use the same model (\texttt{model\_Bacteremia.py}'s
\texttt{dual\_reservoir\_model}), the same treatment protocol, the same
initial condition, and the same fixed parameters (Table~\ref{tab:base}).

\subsection{Role of $EC_{50,L}$ in the model}

$EC_{50,L}$ independently controls how much linezolid effect is achieved at sub-saturating drug concentrations.
\begin{equation}
E_L(t) = \frac{E_{max,L}\, L(t)}{EC_{50,L} + L(t)},
\label{eq:hill}
\end{equation}
the value of $EC_{50,L}$ for each patient will affect how they respond clinically given the identical serum drug concentrations.

\subsection{Protocol and initial condition}

Every simulation in this document integrates the model from the fixed
initial condition
$\mathbf{Y}_0 = (S_b, R_b, S_{res}, R_{res}) = (0,\,0,\,100,\,100)$
CFU$\cdot$mL$^{-1}$ over $t_{\mathrm{end}} = 1944$~h ($81$~d): $21$~d of
untreated growth, $4$~d of vancomycin beginning at $t = 504$~h, $42$~d of
linezolid, and $14$~d of post-treatment follow-up. The limit of detection
is $\mathrm{LOD} = 10$~CFU$\cdot$mL$^{-1}$.

\begin{table}[htbp]
\centering
\begin{tabular}{@{}lll@{}}
\toprule
Parameter & Value & Description \\
\midrule
$\rho_S$              & $0.72$ h$^{-1}$          & Sensitive growth rate, blood \\
$\rho_R$              & $0.649$ h$^{-1}$         & Resistant growth rate, blood ($\sim$10\% fitness cost) \\
$\rho_{res,S}$        & $0.19415$ h$^{-1}$       & Sensitive growth rate, reservoir \\
$\rho_{res,R}$        & $0.1765$ h$^{-1}$        & Resistant growth rate, reservoir ($\sim$10\% fitness cost) \\
$E_{max,V}$           & $1.7$                    & Maximal vancomycin effect \\
$EC_{50,V}$           & $0.25$ mg$\cdot$L$^{-1}$ & Vancomycin half-maximal concentration \\
$E_{max,L}$           & $0.8$                    & Maximal linezolid effect \\
$B_{max,blood}$       & $6\times10^{3}$ CFU$\cdot$mL$^{-1}$ & Blood carrying capacity (shared by $S_b + R_b$) \\
$B_{max,res}$         & $1\times10^{4}$ CFU$\cdot$mL$^{-1}$ & Reservoir carrying capacity (shared by $S_{res} + R_{res}$) \\
$\phi_V$              & $0.15$                   & Vancomycin reservoir penetration fraction \\
$\phi_L$              & $0.30$                   & Linezolid reservoir penetration fraction \\
$f_{r\to b}$          & $5\times10^{-5}$ h$^{-1}$ & Reservoir-to-blood transfer rate \\
$f_{b\to r}$          & $1\times10^{-5}$ h$^{-1}$ & Blood-to-reservoir transfer rate \\
$\mathrm{LOD}$        & $10$ CFU$\cdot$mL$^{-1}$  & Limit of detection \\
\bottomrule
\end{tabular}
\caption{Baseline parameters, identical in both analyses. $EC_{50,L}$ is
absent because it is the quantity under study: swept in
\S\ref{sec:deterministic}, sampled in \S\ref{sec:montecarlo}. Its fixed
project-wide value, used by every \emph{other} script, is
$0.6$~mg$\cdot$L$^{-1}$ --- the same number used here as the Monte Carlo
sampling mean.}
\label{tab:base}
\end{table}


% =============================================================================
\section{Part I --- Deterministic escape threshold}
\label{sec:deterministic}

Every value comes from fixed-parameter integrations, with no statistical sampling.

\subsection{Definition of the threshold}

$EC_{50,L}$ is treated as a bifurcation parameter of the dual-compartment
(blood--reservoir) model, and the resistance-escape threshold
$EC_{50,L}^{*}$ is defined as the value separating trajectories in which
the resistant subpopulation is eradicated by end of follow-up from those in
which it persists.

Persistence is scored on the terminal state of the simulation and not
on transient feature of the trajectory. Let
$R_{res}(t;\,EC_{50,L})$ and $R_b(t;\,EC_{50,L})$ denote the resistant
densities in the reservoir and blood compartments. Defining the residual
end-of-simulation burden relative to the limit of detection,
\begin{equation}
  f_c(EC_{50,L}) \;=\; X_c\!\left(t_{\mathrm{end}};\,EC_{50,L}\right) - \mathrm{LOD},
  \qquad c \in \{b,\,res\},
  \label{eq:objective}
\end{equation}
where $X_b \equiv R_b$ and $X_{res} \equiv R_{res}$, the threshold is the
root
\begin{equation}
  f_c\!\left(EC_{50,L}^{*}\right) = 0 .
  \label{eq:root}
\end{equation}
Values $EC_{50,L} > EC_{50,L}^{*}$ (reduced potency, since a larger
$EC_{50,L}$ requires a higher drug concentration for the same effect via
Eq.~\eqref{eq:hill}) yield $f_c > 0$ and are classified as escape;
values below yield $f_c < 0$ and are classified as suppression. This approach provides an explicit and quantitative criterion for separating clinical response from relapse.

\subsection{Numerical procedure}

The threshold is located in two stages.

\paragraph{Stage 1: deterministic parameter sweep.}
A single simulation is run at each of $121$ values of $EC_{50,L}$ on a
uniform grid spanning $[0.1,\,1.3]$~mg$\cdot$L$^{-1}$ (step
$0.01$~mg$\cdot$L$^{-1}$), with every other parameter, the initial
condition, the dosing schedule, and the pharmacokinetic concentration
profiles held fixed, so that $EC_{50,L}$ is the sole source of variation
between runs. Only the terminal values $R_b(t_{\mathrm{end}})$ and
$R_{res}(t_{\mathrm{end}})$ are retained. This sweep characterizes the
shape of the dose response curve.  It confirms that
the transition is a single, monotone switch and not a gradual
decline or a multi-valued response.  These procedure ensures there is only one root.


\paragraph{Stage 2: root-finding.}
The precise crossing is obtained by solving Eq.~\eqref{eq:root} with
Brent's method \cite{brent1973}, a bracketing algorithm combining
bisection with inverse quadratic interpolation. Each evaluation of $f_c$
entails one complete ODE integration over the full protocol. The root is
bracketed by the sweep endpoints as noted above and is solved to an absolute tolerance of
$10^{-3}$~mg$\cdot$L$^{-1}$. Prior to bisection, the sign condition
\begin{equation}
  f_c(0.1) < 0 \quad\text{and}\quad f_c(1.3) > 0
  \label{eq:bracket}
\end{equation}
is verified explicitly (failure of either condition indicates that no sign
change and no threshold within the interval).
Equations~\eqref{eq:objective}--\eqref{eq:bracket} are solved independently
for the blood ($c = b$) and reservoir ($c = res$) compartments, yielding
two separate thresholds.

\subsection{Numerical implementation}

All integrations use the LSODA solver as exposed by
\texttt{scipy.integrate.odeint}, with relative and absolute tolerances of
$10^{-7}$ and $10^{-9}$ and a maximum of $5000$ internal steps per output
interval; solutions are evaluated on a fixed grid of $1150$ time points
spanning $[0,\,t_{\mathrm{end}}]$. Root-finding uses
\texttt{scipy.optimize.brentq}. The solver tolerances are several orders of
magnitude tighter than the LOD-scale quantities compared in
Eq.~\eqref{eq:objective}, ensuring that the located threshold reflects
model dynamics rather than integration error.

\subsection{Result}
\label{sec:threshold_result}

Under Table~\ref{tab:base}, the resistant reservoir and the resistant blood
compartment both cross the limit of detection at
\begin{equation}
EC_{50,L}^{*} = 0.488~\mathrm{mg}\cdot\mathrm{L}^{-1},
\end{equation}
the two crossings agreeing to within the $10^{-3}$~mg$\cdot$L$^{-1}$
bisection tolerance and therefore not resolved as separate values. They are
expected to coincide because both experience the same linezolid effect term
through the shared blood-compartment efficacy and are coupled by
bidirectional bacterial exchange. The reservoir value is adopted as the
operative threshold in downstream analyses, since it is the reservoir that
seeds post-treatment relapse.

The threshold is re-derived by root-finding it is recomputed when upstream parameters are revised.
The following are examples:

\begin{itemize}
\item raising $\rho_S$ from $0.60$ to $0.61$~h$^{-1}$ moved the threshold
from $0.988$ to $0.932$~mg$\cdot$L$^{-1}$;
\item granting the sensitive strain a $10\%$ reservoir growth advantage
($\rho_{res,S} = 0.19415$ versus $\rho_{res,R} = 0.1765$~h$^{-1}$), with a
compensating reduction of $\phi_L$ from $0.45$ to $0.30$, moved it to
$0.868$~mg$\cdot$L$^{-1}$;
\item a subsequent rise in $\rho_S$ from $0.61$ to $0.72$~h$^{-1}$ moved it
to its present $0.488$~mg$\cdot$L$^{-1}$.
\end{itemize}

\subsection{Assumptions and limitations}

Two assumptions underlie the procedure. First, Brent's method returns a
single root within the bracketing interval, so the threshold is well
defined only if the terminal resistant burden is monotone in $EC_{50,L}$
over $[0.1,\,1.3]$~mg$\cdot$L$^{-1}$. The $121$-point grid sweep of
Stage~1 is the check on this assumption, and confirms a single sharp
transition under the present parameterisation; a non-monotone response
arising from a future re-parameterisation would render the reported root
non-unique. Second, the bracketing interval is fixed a priori, and a
parameter revision displacing the threshold outside
$[0.1,\,1.3]$~mg$\cdot$L$^{-1}$ would violate Eq.~\eqref{eq:bracket}; that
condition is tested explicitly at each run, so such a displacement is
reported as an absent threshold rather than propagating silently as an
incorrect value.


% =============================================================================
\section{Part II --- Monte Carlo patient population}
\label{sec:montecarlo}

\subsection{Sampling distribution}

$EC_{50,L}$ is sampled from a mean-corrected log-normal distribution,
\begin{equation}
\sigma_{\ln} = \sqrt{\ln(1+CV^2)}, \qquad
\mu_{\ln} = \ln(EC_{50,L,\mathrm{base}}) - \tfrac{1}{2}\sigma_{\ln}^2, \qquad
EC_{50,L} \sim \mathrm{LogNormal}(\mu_{\ln}, \sigma_{\ln}^2),
\label{eq:lognormal}
\end{equation}
with $EC_{50,L,\mathrm{base}} = 0.6$~mg/L and $CV = 1.1$, so that
$\mathbb{E}[EC_{50,L}]$ lands exactly on baseline.  The same
mean-corrected, CV-parameterized convention is used by
\texttt{MC\_immune\_response.py} and \texttt{MC\_emax\_lzd.py}. The
$EC_{50,L}$ is a single value used throughout the project.  Numerically:
$\sigma_{\ln} = 0.8905$ and $\mu_{\ln} = -0.9073$.

\subsection{Procedure, as implemented}

The code's order of execution.

\begin{enumerate}
\item \textbf{Load the model.} \texttt{model\_Bacteremia.py} is loaded
dynamically via \texttt{importlib} (rather than a plain \texttt{import}) so
a missing file produces an explicit error rather than an
\texttt{ImportError} traceback.

\item \textbf{Fix run-level constants.} \texttt{NUM\_ITERATIONS = 1000},
\texttt{LOD = 10.0} CFU/mL, \texttt{CV = 1.1}, \texttt{SEED = 44},
\texttt{total\_h = 1944}~h, \texttt{vanco\_start = 504}~h, and a shared
evaluation grid \texttt{t\_eval = linspace(0, 1944, 1150)} used by every
ODE solve.

\item \textbf{Build the drug-concentration and immune-response callables
once.} \texttt{PharmacokineticModel().concentration\_function(...)} is
called once each for vancomycin and linezolid, and
\texttt{ImmuneResponse()} is instantiated once. These closures are reused
unmodified by all $1000$ solves --- only $EC_{50,L}$ changes between
iterations.

\item \textbf{Fix \texttt{BASE\_PARAMS} and the initial state}
(Table~\ref{tab:base}, $\mathbf{Y}_0$ above). $EC_{50,L}$ is 
absent from \texttt{BASE\_PARAMS}.  It is merged in per call at step~7.

\item \textbf{Compute the log-normal parameters} of
Eq.~\eqref{eq:lognormal} once.

\item \textbf{Draw all $1000$ samples at once.} A seeded generator,
\texttt{rng = np.random.default\_rng(44)}, draws the full vector in a
single call.  The run is fully reproducible and sample $i$ is used by
iteration $i$.

\item \textbf{Compute the internal deterministic switch curve
($131$ points)} over $[0.05,\,3.3]$~mg/L, recording final $R_b$ and
$R_{res}$ at each point, then root-find \texttt{ESCAPE\_THRESH} on the same
objective as Eq.~\eqref{eq:root}. If the endpoint signs fail to bracket a
root, the code falls back to the coarse grid crossing.

\item \textbf{Main Monte Carlo loop ($1000$ iterations).} Pre-allocated
\texttt{NaN}-filled arrays (each $1000 \times 1150$) are filled by, for
each $i$: building \texttt{params\_i} as a shallow copy of
\texttt{BASE\_PARAMS} with only $EC_{50,L}$ replaced by sample $i$;
integrating from the same $\mathbf{Y}_0$ over the same \texttt{t\_eval};
flooring each of the four trajectories to $0$ wherever it falls below the
LOD. This ensures that sub-detection noise reads as true clearance and not a small positive count.  The full trajectory is stored into row $i$. All run-to-run
variance is attributable to sampled $EC_{50,L}$ alone. 

\item \textbf{Classify each patient.} \texttt{final\_R\_res} is the last
time point of each row, and \texttt{escaped\_mask = final\_R\_res > LOD}
marks each patient as escape or suppression by their \emph{own} simulated
outcome.

\item \textbf{Render the two figures} (\S\ref{sec:outputs}) directly from
the stored arrays.  No further simulation is performed.

\item \textbf{Print a console summary} of the sampled $EC_{50,L}$
distribution and, where any samples exceed the LOD, the median and
5th/95th percentiles of peak $S_b$ and peak $R_b$.
\end{enumerate}

\subsection{Result}

Against the $0.488$~mg/L threshold, the sampled population splits
\begin{center}
\textbf{$43.6\%$ escape ($436/1000$)} \quad versus \quad
\textbf{$56.4\%$ suppressed ($564/1000$)},
\end{center}
with sampled median $0.43$~mg/L (5th--95th percentile
$0.10$--$1.68$~mg/L).

This fraction is largely a readout of how close the calibration sits to its
own switch. Applying Eq.~\eqref{eq:lognormal}, the sampling mean
($0.6$~mg/L) and the threshold ($0.488$~mg/L) are separated by
\begin{equation}
z = \frac{\ln(0.488) - \mu_{\ln}}{\sigma_{\ln}}
  = \frac{\ln(0.488) - (-0.9073)}{0.8905} \approx 0.213,
\end{equation}
about a fifth of a standard deviation of $\ln X$. The standard-normal upper
tail at that $z$ predicts $1 - \Phi(0.213) = 0.416$, i.e.\ $41.6\%$ escape,
within two points of the simulated $43.6\%$.  It illustrates that a calibration that close to its own switch necessarily splits the population close to evenly.

\subsection{Why the two stages, in this order}
\label{sec:agreement}

The deterministic sweep answers \emph{where the switch is}; the Monte Carlo
loop answers \emph{how a population distributes around it}. Keeping them
separate matters for one specific reason: in the figures, each Monte Carlo
sample is coloured by its \emph{own} simulated outcome (whether that
patient's final $R_{res}$ exceeds the LOD), not by which side of
\texttt{ESCAPE\_THRESH} its sampled $EC_{50,L}$ happens to fall on. Were
the classification taken from the threshold, the population panel would
agree with the switch panel by construction and could demonstrate nothing.
As built, the deterministic switch is a prediction and the Monte Carlo
colouring is the outcome, so any disagreement is visible rather than
hidden.

The same independence applies across scripts. \texttt{MC\_ec50\_lzd.py}
bisects over $[0.05,\,3.3]$ on a $131$-point grid, while
\texttt{EC50\_lzd\_threshold\_sweep.py} bisects over $[0.1,\,1.3]$ on a
$121$-point grid. The intervals no longer share an endpoint, and the grids
differ in both extent and spacing, yet both locate the threshold at
$0.488$~mg/L. The wider Monte Carlo bracket is required because that script
must cover its own sampled population: at mean $0.6$ with $CV = 1.1$, the
1st--99th percentile of the log-normal draw spans roughly
$[0.05,\,3.20]$~mg/L, which the narrower diagnostic interval would not
contain.


% =============================================================================
\section{Figures}
\label{sec:outputs}

Four figures cover this analysis. Two are purely deterministic, one is
purely Monte Carlo, and one is a combination of both methods.  This is an important distinction, because the deterministic scatter plot superficially does resembles a Monte Carlo.

\paragraph{\texttt{ec50\_lzd\_threshold\_sweep.png} --- DETERMINISTIC.}
Final $R_b$ and $R_{res}$ across the $121$-point grid, escape zones shaded,
both bisected thresholds marked. Sub-LOD values are floored to
$\mathrm{LOD}/2$ for display so the curves stay continuous and the
transition reads as a visible jump rather than a gap.

\paragraph{\texttt{ec50\_terminal\_burden\_vs\_ec50.png} --- DETERMINISTIC.}
The same runs as a categorical outcome plot: $24$ markers subsampled from
the same grid, $R_{res}$ only. Open circles are pinned at the LOD for runs
that clear (their true residual is $\sim$$10^{-12}$ and would fall off a
log axis); filled circles sit at the true final count for runs that escape.
\emph{Each marker is one fixed-parameter run.}

\paragraph{\texttt{mc\_ec50\_lzd\_sweep\_kinetics.png} --- MONTE CARLO.}
A four-panel time course ($S_b$, $R_b$, $S_{res}$, $R_{res}$) over the
$1944$-h window, each panel showing $20$ individual sample trajectories
plus a median and $5^{\mathrm{th}}$--$95^{\mathrm{th}}$-percentile band,
with the dosing windows overlaid and all four panels sharing a
$0$--$81$~day axis.

The median and band are taken over the \emph{escaped} patients only
($n = 436$), not all $1000$ iterations. Sub-LOD values are stored as $0$,
so an all-patient median is itself $0$ wherever escape (clinical relapse) is a minority outcome making the relapse disappear from the plot. Restricting to escaped rows and
using \texttt{nanmedian} keeps it visible. A gap therefore means ``not
detectable in any escaped patient'', not ``no data'', and the band is a
conditional quantity: the spread among patients who escaped.

\paragraph{\texttt{mc\_ec50\_lzd\_switch.png} --- MIXED, one source per
panel.} The top panel is the continuous \emph{deterministic} switch curve
(final $R_b$/$R_{res}$ versus $EC_{50,L}$ from the $131$-point sweep) with
\texttt{ESCAPE\_THRESH} marked and the suppressed/escape regimes shaded.
The bottom panel is a \emph{Monte Carlo} histogram of the $1000$ sampled
$EC_{50,L}$ values, split and coloured by each patient's actual simulated
outcome. Both panels share an $x$-axis spanning the sweep interval, which
clips display of the roughly top $1\%$ of draws above $3.3$~mg/L; those
samples are still counted in the reported escape fraction.

\begin{thebibliography}{9}
\bibitem{brent1973}
Brent, R.~P. (1973). \emph{Algorithms for Minimization without
Derivatives}, ch.~4. Prentice-Hall, Englewood Cliffs, NJ.

\bibitem{brent1971}
Brent, R.~P. (1971). ``An algorithm with guaranteed convergence for
finding a zero of a function''. \emph{Computer Journal} 14, 422--425.
\end{thebibliography}

\end{document}
